\documentclass[11pt]{article}

\usepackage[letterpaper,margin=1in]{geometry}
\usepackage{amsmath,amssymb}
\usepackage{booktabs}
\usepackage{longtable}
\usepackage{array}
\usepackage{xcolor}
\usepackage{listings}
\usepackage{url}
\usepackage{hyperref}
\usepackage{titlesec}
\usepackage{parskip}
\usepackage{fancyhdr}
\usepackage{enumitem}

\hypersetup{
  colorlinks=true,
  linkcolor=black,
  urlcolor=blue!50!black,
  citecolor=black,
  pdftitle={KV Cache Engine --- Reference Model API},
  pdfauthor={LonghornSilicon},
  pdfsubject={API Reference, Version 0.2}
}

\titleformat{\section}{\normalfont\Large\bfseries}{\thesection.}{0.6em}{}
\titleformat{\subsection}{\normalfont\large\bfseries}{\thesubsection}{0.6em}{}
\titleformat{\subsubsection}{\normalfont\normalsize\bfseries}{\thesubsubsection}{0.6em}{}

\pagestyle{fancy}
\fancyhf{}
\lhead{\small KV Cache Engine --- Reference Model API}
\rhead{\small v0.2}
\cfoot{\thepage}
\renewcommand{\headrulewidth}{0.4pt}

\definecolor{codebg}{gray}{0.97}
\definecolor{codekey}{rgb}{0.13,0.29,0.53}
\definecolor{codecom}{rgb}{0.25,0.5,0.25}
\lstset{
  basicstyle=\ttfamily\small,
  backgroundcolor=\color{codebg},
  frame=single,
  framerule=0.4pt,
  rulecolor=\color{black!30},
  xleftmargin=0.5em,
  xrightmargin=0.5em,
  aboveskip=0.6em,
  belowskip=0.6em,
  showstringspaces=false,
  breaklines=true,
  keywordstyle=\color{codekey}\bfseries,
  commentstyle=\color{codecom}\itshape,
  numbers=none
}

\setlist{itemsep=1pt,topsep=2pt}

\begin{document}

%==============================================================================
\thispagestyle{empty}
\begin{center}
\vspace*{1.2in}
{\Huge\bfseries KV Cache Engine}\\[0.4em]
{\LARGE\bfseries Reference Model API}\\[1.2em]

{\large \textsc{LonghornSilicon LLM Inference Accelerator}}\\[0.4em]
{\large Block 2 of 4 --- Bit-Accurate C++ / Python Models (ChannelQuant)}\\[2em]

\rule{0.7\textwidth}{0.4pt}\\[1em]

\begin{tabular}{rl}
\textbf{Version}:    & 0.2 (ChannelQuant codec) \\
\textbf{Date}:       & 2026-07-03 \\
\textbf{Status}:     & Bit-exact vs golden vectors + RTL (3-way parity) \\
\textbf{Source}:     & \texttt{sw/reference\_model/channelquant\_ref.\{hpp,cpp\}} \\
\end{tabular}\\[1em]
\rule{0.7\textwidth}{0.4pt}
\vfill
\end{center}
\clearpage

%==============================================================================
\tableofcontents
\clearpage

%==============================================================================
\section{Overview}
\label{sec:overview}

The reference model is a bit-accurate implementation of the \textbf{ChannelQuant}
KV-cache codec (the RTL's behavioral oracle, ported 1:1). It is a small
free-function API in namespace \texttt{lhsi::cq}: compress / decompress a single
KV head of shape $[T, D]$ (row-major fp16) through the per-channel key path and
the per-token value path.

\begin{table}[h]
\centering
\small
\begin{tabular}{@{}l l l l@{}}
\toprule
Codec & Namespace & Entry points & Tests \\
\midrule
\textbf{ChannelQuant} & \texttt{lhsi::cq} & \texttt{compress\_keys/values}, \texttt{decompress\_*} & 9 golden vectors (C++), 9 (Python) \\
\bottomrule
\end{tabular}
\end{table}

\noindent The Python algorithm reference and golden vectors are frozen in the
sibling repository \path{../channelquant/reference/}. The retired TurboQuant+ C++
reference (\path{kv_cache_engine_ref.{hpp,cpp}}) is retained in the same directory
for archival/regression only; it is not the shipped codec. Top-level orientation:
\path{sw/README.md}. ISA: \path{docs/isa/kv_cache_engine_isa.pdf}.

%==============================================================================
\section{Building and Testing}
\label{sec:build}

\begin{lstlisting}[language=bash]
make test-cq       # ChannelQuant C++ tests vs the 9 golden vectors
make test-py       # Python reference tests
make test-all      # legacy TQ + ChannelQuant + Python
make shared        # libkv_cache_engine_ref.so
make static        # libkv_cache_engine_ref.a
\end{lstlisting}

\noindent Requires Python 3.10+, NumPy, and a C++17 compiler. The tests load the
vendored golden vectors under
\path{rtl/tb/testvectors/channelquant/} (committed; no generation step).

%==============================================================================
\section{ChannelQuant Reference API}
\label{sec:cqapi}

\subsection{Compress / Decompress}

Single KV head, shape $[T, D]$ row-major fp16 (as \texttt{uint16\_t} bit
patterns). \texttt{bits} = 8 (CQ-8) or 4 (CQ-4/CQ-4+).

\begin{lstlisting}[language=C++]
#include "channelquant_ref.hpp"
using namespace lhsi::cq;

// Values: per-token scale -> INT4/INT8.
ValueBlob v = compress_values(V_f16, T, D, /*bits=*/4);
std::vector<uint32_t> Vhat_f32 = decompress_values(v);   // [T*D] fp32 bits

// Keys: per-channel scale over groups of G tokens -> INT4; the channels in
// `outlier_idx` are held FP16 (CQ-4+). Empty outlier_idx = plain CQ-4.
KeyBlob k = compress_keys(K_f16, T, D, /*bits=*/4, /*G=*/128, outlier_idx);
std::vector<uint32_t> Khat_f32 = decompress_keys(k);     // [T*D] fp32 bits
\end{lstlisting}

\subsection{Core Cores (one per RTL unit)}

\begin{lstlisting}[language=C++]
uint16_t s   = scale_from_amax(amax_f16, bits);   // cq_scale_unit: max(amax/qmax, EPS)
int      q   = quant_code(x_f16, s, bits);        // cq_quant: round-half-even + clamp
uint32_t xh  = dequant_f32(q, s);                 // cq_dequant: code*scale -> fp32
std::vector<uint8_t> p4 = pack_int4(codes);       // contract-ordered nibble packing
std::vector<uint8_t> p8 = pack_int8(codes);
// fp helpers: f16_to_real / f32_to_real / real_to_f16 / real_to_f32
\end{lstlisting}

%==============================================================================
\section{Numerical Semantics (Frozen)}
\label{sec:semantics}

All arithmetic is fixed-function fp16/fp32, matching the RTL exactly
(contract v0.2):

\begin{itemize}
\item \textbf{Elements}: IEEE-754 fp16 inputs; fp32 decompressed outputs
      (read bus, contract \S1).
\item \textbf{Scale}: $s = \max(\text{amax}/q_{\max},\ \text{EPS})$ cast to fp16,
      with $\text{EPS}=2^{-14}$ and $q_{\max}=7$ (INT4) or $127$ (INT8).
\item \textbf{Quantize}: $\mathrm{clamp}(\mathrm{round\_half\_even}(x/s))$ to the
      tier's signed range (INT4 $[-8,7]$, INT8 $[-128,127]$).
\item \textbf{Dequantize}: $\hat{x} = \text{code}\cdot s$ computed exactly in fp32
      (the product fits the 23-bit mantissa, no rounding).
\item \textbf{Key axis}: per channel over a group of $G$ tokens. \textbf{Value
      axis}: per token. \textbf{Outlier} channels (CQ-4+): stored raw fp16.
\end{itemize}

\noindent There is no rotation, no Johnson--Lindenstrauss projection, and no
Lloyd--Max codebook (those were the retired TurboQuant+ codec). If the model
disagrees with the chip on any input, the model is wrong --- open an issue with
the failing vector.

%==============================================================================
\section{Compressed Data Structures}
\label{sec:structures}

\subsection{\texttt{ValueBlob} (per-token path)}

\begin{lstlisting}[language=C++]
struct ValueBlob {
    int T, D, bits;
    std::vector<uint16_t> scales;   // [T] fp16 per-token scales
    std::vector<int>      codes;    // [T*D] signed codes (row-major)
    std::vector<uint8_t>  payload;  // packed byte stream (contract packing)
};
\end{lstlisting}

\subsection{\texttt{KeyBlob} (per-channel grouped path)}

\begin{lstlisting}[language=C++]
struct KeyBlob {
    int T, D, bits, G;
    std::vector<std::pair<int,int>> groups;  // [(a,b), ...] token ranges
    std::vector<int>      keep;              // non-outlier channel indices
    std::vector<int>      outlier;           // outlier channel indices (sorted)
    std::vector<uint16_t> scales;            // per group, nk fp16 (concatenated)
    std::vector<uint8_t>  payload;           // per group, packed INT4 (concatenated)
    std::vector<uint16_t> sidecar;           // [T*k] fp16, t-major (outlier cols)
};
\end{lstlisting}

\noindent The RTL folds this into a single per-channel SRAM record (keep channel
$\to\{$scale, INT4$\}$; outlier channel $\to\{$fp16, code $+1\}$); the blob keeps
the fields separate so each surface (scales / payload / K̂ / sidecar) can be
checked independently against the golden vectors.

%==============================================================================
\section{Verification Status}
\label{sec:verification}

\begin{table}[h]
\centering
\small
\begin{tabular}{@{}l l l@{}}
\toprule
Test & Aspect & Status \\
\midrule
9 golden vectors (C++) & scales / payload / K̂ / V̂ / sidecar & Bit-exact \\
9 golden vectors (Python) & same surfaces & Bit-exact \\
3-way Python↔C++↔SV parity & compress + decompress & Closed (contract \S8) \\
Tiers CQ-8 / CQ-4 / CQ-4+ & $D\in\{64,128\}$, full + partial groups & Pass \\
RTL bit-exact vs golden & \texttt{sim\_cq} / \texttt{sim\_kpath} / \texttt{sim\_top} & Pass \\
\bottomrule
\end{tabular}
\end{table}

\noindent Because the C++ core is a verbatim port of the SV core, C++$\equiv$SV by
construction, not just coincidence on these vectors.

%==============================================================================
\section{Co-Design Context}
\label{sec:codesign}

This directory is the deliverable for compiler-team integration: the reference
model exposes the same numerical contract the hardware implements, so compiler
backends can target the KV cache engine before silicon or FPGA prototypes exist.

\medskip
\noindent Next integration steps: (1) compiler targets the Python/C++ reference
(now); (2) AXI-Lite + AXI-Stream on Zynq UltraScale+ FPGA; (3) multi-block FPGA
project with the precision controller; (4) TSMC 16FFC silicon (2027+).

\vspace{0.5in}
\noindent\rule{\textwidth}{0.4pt}\\
\noindent\small\textit{LonghornSilicon KV Cache Engine --- Reference Model API.
This document is open and may be redistributed.}

\end{document}
